% Part: incompleteness
% Chapter: introduction
% Section: historical-background

\documentclass[../../../include/open-logic-section]{subfiles}

\begin{document}

\olfileid{inc}{int}{bgr}

\olsection{تاريخي شاليد}

په دې برخه کښې به پر هغو تاريخي پرمختګونو لنډ بحث وکړو چې د
ناتکميلۍ د قضيو د چاپېريال په پوهېدو کښې مرسته کوي۔ په ځانګړي ډول به
د رياضياتي منطق د تاريخ يوه ډېره لنډه ټوليزه کتنه وړاندې کړو، او بيا
به د رياضياتو د بنسټونو د تاريخ په باب څو خبرې وکړو۔

\begin{digress}
  د «رياضياتي منطق» غونډله دوه‌معناييزه ده۔ د «رياضياتي» ويے د موضوع
  د بيان په توګه تعبيرېدے شي، لکه «د رياضياتو منطق»، چې د رياضياتي
  استدلال اصول ښيي؛ يا د لارو د بيان په توګه، لکه «د منطق رياضيات»،
  چې د استدلال د اصولو رياضياتي مطالعه ښيي۔ لاندې روايت رياضياتي منطق
  په دواړو معناوو، او ډېر ځله په يوه وخت کښې، رانغاړي۔
\end{digress}

د منطق مطالعه په اصل کښې له ارسطو څخه پيل شوه، چې نږدې په
384--322 \textsc{ق‌م} کښې ژوندے و۔ د هغه \emph{مقولات}،
\emph{لومړۍ تحليلات} او \emph{دويمې تحليلات} د علمي استدلال د اصولو
منظمې مطالعې رانغاړي، چې د قياس بشپړه او منظمه مطالعه هم پکې شامله ده۔

د ارسطو منطق د منځنيو پېړيو په اوږدو کښې پر مدرسي فلسفې واکمن و؛ ان
تر اتلسمې پېړۍ پورې کانټ ټينګار کاوه چې د ارسطو منطق بشپړ دے او
سمون ته اړتيا نۀ لري۔ خو د قياس تيوري ډېره محدوده ده او د رياضياتي
استدلال له تر ټولو سطحي اړخونو پرته نور څه نۀ شي مدل کولے۔ يوه پېړۍ
تر دې وړاندې، د نيوټن هم‌مهالي لايبنيڅ د منطقي استدلال دپاره د يوه
بشپړ «حساب» تصور وکړ او د داسې حساب د جوړولو پر لور ئې ځينې لومړني
ګامونه واخيستل؛ په اصل کښې ئې د قضييز منطق يوه بڼه بيان کړه۔

نولسمه پېړۍ د منطق دپاره يو بنسټيز بدلون و۔ په 1854 کښې جورج بول
\emph{د فکر قانونونه} وليکل، چې د قضييز منطق بشپړه الجبري مطالعه
پکې وه او له اوسنيو وړاندېونو څخه ډېره لرې نۀ ده۔ په 1879 کښې ګوټلوب
فرېګې خپله \emph{بېګريفشريفټ} (مفهومي ليکدود) خپره کړه، چې قضييز
منطق د کميت‌ټاکونکو او اړيکو په ورزياتولو غځوي او په دې توګه د لومړۍ
درجې منطق هم رانغاړي۔ په حقيقت کښې د فرېګې منطقي نظامونو د لوړې درجې
منطق او نور څه هم رانغاړل۔ فرېګې په خپلو \emph{د حساب بنسټيزو
قانونونو} کښې هڅه وکړه وښيي چې ټول حساب د هغه په بېګريفشريفټ کښې له
سوچه منطقي فرضيو څخه استنتاجېدے شي۔ له بده مرغه دا فرضيې ناسازګارې
وختې، لکه رسل چې په 1902 کښې وښودل۔ خو که ناسازګار بديهي اصل څنګ ته
کړو، فرېګې نږدې په يوازې ځان اوسنی منطق اختراع کړ، چې يوه حيرانوونکې
لاسته راوړنه ده۔ له بول وروسته الجبري فکر لرونکو پوهانو، لکه پيرس او
شروډر، کميتي منطق په خپلواکه توګه هم جوړ کړ۔

اوس راځئ د رياضياتو د بنسټونو پرمختګونو ته مخه کړو۔ البته څرنګه چې
منطق په رياضياتو کښې مهم رول لري، له پاس بيان شويو پرمختګونو سره ئې
ډېر متقابل اغېز شته۔ د بېلګې په توګه، فرېګې خپل منطق په دې څرګنده
موخه جوړ کړ چې وښيي ټول رياضيات يوازې پر هغه منطقي چوکاټ ولاړېدے شي؛
په ځانګړي ډول غوښتل ئې وښيي چې رياضيات، د کانټ د نظر پر خلاف، له
تجربې وړاندې \emph{تحليلي} حقايق دي، نه له تجربې وړاندې
\emph{ترکيبي} حقايق۔

ډېر خلک د رياضياتو اصلي زېږون له يونانيانو سره تړي۔ د اقليدس
\emph{اصول}، چې نږدې 300 \textsc{ق‌م} کښې ليکل شوي، لا له وړاندې د
يوناني رياضياتو يوه پخه بېلګه ده او پر کره‌والي او دقت ټينګار کوي۔
د اقليدس د \emph{اصولو} تعريفونه او ثبوتونه د نن ورځې د لېسې د هندسې
په درسي کتابونو کښې نږدې هماغسې پاتې دي (تر هغه ځايه چې لا په لېسو
کښې هندسه تدريسېږي)۔ د رياضياتي استدلال دا مدل نه يوازې په رياضياتو
بلکې د فلسفې په څانګو کښې هم د کره استدلال بېلګه ګڼل شوې ده۔
(سپينوزا ان اخلاقي او ديني استدلالونه د اقليدسي سبک له مخې وړاندې
کړل، چې ليدل ئې عجيبه دي!)

نيوټن او لايبنيڅ په اوولسمه پېړۍ کښې کالکولس اختراع کړ۔ (د لومړيتوب
يوه سخته شخړه پېړۍ پېړۍ روانه وه، خو اوس زياتره پوهان مني چې دواړه
پرمختګونه تر ډېره خپلواک وو۔) په کالکولس کښې، د بېلګې په توګه، د
نامتناهي وړو مقدارونو د نامتناهي جمعو په باب استدلال راځي؛ دغو
ځانګړتياوو د بيشپ برکلي نيوکې راوپارولې، چې استدلال ئې کاوه پر خدای
باور د هغه د وخت تر رياضياتو لږ معقول نۀ دے۔ په اتلسمه پېړۍ کښې د
کالکولس لارې په پراخه توګه وکارول شوې؛ مثلاً ليونارډ اويلر د
نامتناهي جمعو محاسبې په حيرانوونکو پايلو وکارولې۔

په نولسمه پېړۍ کښې رياضيپوهانو هڅه وکړه د کالکولس د لا ټينګ بنسټ په
جوړولو د برکلي نيوکو ته ځواب ووايي۔ د کوشي، وايرشتراس، بولتسانو او
نورو هڅو موږ د حد، پيوستون، تفاضل او تکامل اوسنيو تعريفونو ته
ورسولو، چې د «اېپسيلونونو او ډېلټاګانو» په اصطلاحاتو کښې دي؛ يعنې
نامتناهي وړو مقدارونو ته هېڅ مراجعه نۀ کوي۔ د پېړۍ په وروستيو کښې
رياضيپوهانو هڅه وکړه لا وړاندې لاړ شي او د حقيقي عددونو په ګډون د
کالکولس ټول اړخونه د طبيعي عددونو په اصطلاحاتو کښې بيان کړي۔
(کرونېکر: «خدای صحيح عددونه پيدا کړل، نور هر څه د انسان کار دے۔»)
په 1872 کښې ډېدېکينډ «پيوستون او غيرناطق عددونه» وليکل، چې وئې ښودل
څنګه حقيقي عددونه د ناطقو عددونو د سټونو په توګه «جوړېدے» شي (او
ناطق عددونه، لکه چې پوهېږئ، د طبيعي عددونو د جوړو په توګه ليدل کېدے
شي)؛ په 1888 کښې ئې «Was sind und was sollen die Zahlen» (نږدې معنا
ئې «طبيعي عددونه څه دي، او څه بايد وي؟») وليکل، چې موخه ئې په سوچه
«منطقي» اصطلاحاتو کښې د طبيعي عددونو بيان و۔ په 1887 کښې کرونېکر
«\"Uber den Zahlbegriff» («د عدد د مفهوم په باب») وليکل او د صحيح
عددونو په اصطلاحاتو کښې ئې د ټولو رياضياتي څيزونو د تمثيل خبره وکړه؛
په 1889 کښې جوزېپې پيانو د طبيعي عددونو صوري، نښيز بديهي اصول ورکړل۔

د نولسمې پېړۍ پاى د نامتناهي په چلند کښې يو نوی زړورتوب هم راوړ۔
تر دې وړاندې له نامتناهي څيزونو او جوړښتونو (لکه د طبيعي عددونو له
سټ) سره ډېر په احتياط چلند کېده؛ «نامتناهي ډېر» د «څومره چې وغواړئ
هومره» په معنا، او «په حد کښې ورنږدې کېږي» د «څومره چې وغواړئ هومره
نږدې کېږي» په معنا ګڼل کېده۔ خو ګيورګ کانتور وښودل چې نامتناهي په
خپله ظاهري معنا اخيستل کېدے شي۔ د کانتور، ډېدېکينډ او نورو کار د
رياضياتو هغه عمومي سټ‌تيوريکي پوهاوی رامنځته کړ چې اوس په پراخه توګه
منل شوے دے۔

دا مو په منطق او بنسټونو کښې د شلمې پېړۍ پرمختګونو ته رسوي۔ په 1902
کښې رسل د فرېګې په منطقي نظام کښې پاراډوکس وموند۔ په 1904 کښې
څېر مېلو د کانتور د ښه‌ترتيب اصل د تش په نوم «د انتخاب بديهي اصل» په
کارولو ثابت کړ؛ د دې بديهي اصل د رواوالي په باب ډېر بحث وشو۔ د 1910
او 1913 تر منځ د رسل او وايټ‌هېډ د \emph{پرېنسيپيا مېتېماتيکا} درې
ټوکونه خپاره شول، چې د رياضياتو پر منطقي بنسټونو د درولو د فرېګې
پروګرام ئې وغځاوه۔ له بده مرغه رسل او وايټ‌هېډ اړ شول دوه داسې اصول
ومني چې د سوچه منطقي اصولو په توګه ئې توجيه ستونزمنه ښکارېده: د
نامتناهي والي بديهي اصل او د «راکمېدنې» بديهي اصل۔ په 1900مو کلونو
کښې پوانکاره په رياضياتو کښې د «غيرپريديکاتي تعريفونو» پر کارولو
نيوکه وکړه، او په 1910مو کلونو کښې بروور د ټولو رياضياتو د بيا بنسټيز
کولو دپاره د داسې «شهودي» بنسټ وړانديز پيل کړ چې د خارج وسط اصل
($!A \lor \lnot !A$) پکې نۀ کارېده۔

په رښتيا هم عجيب وختونه وو!{} منطق ته د ټولو رياضياتو د راکمولو
پروګرام اوس «منطقپالنه» بلل کېږي او د پاسنيو ستونزو له امله عموماً
ناکام ګڼل کېږي۔ د شهودي ذهني جوړونو په اصطلاحاتو کښې د رياضياتو د
جوړولو پروګرام «شهوديت» بلل کېږي او پر ورځنيو رياضياتو ډېر سخت
محدوديتونه ايښوونکی ګڼل کېږي۔ د پېړۍ د بدلېدو شاوخوا ډېوېډ هېلبرټ،
چې د هر وخت له تر ټولو اغېزمنو رياضيپوهانو څخه و، د کانتور او
ډېدېکينډ د نويو مجردو لارو ټينګ ملاتړی و: «هېڅوک به موږ له هغه جنت
څخه ونه باسي چې کانتور راته جوړ کړے دے۔» په همدې وخت کښې د دې نويو
لارو (چې عجيبه دا ده اوس «کلاسیکې» بلل کېږي) بنسټيزو نيوکو ته هم
متوجه و۔ هغه داسې لار وړانديز کړه چې دواړه موخې تر لاسه کړي:
\begin{enumerate}
\item کلاسیکې لارې په صوري بديهي اصولو او قاعدو تمثيل کړئ؛ رياضياتي
  پوښتنې په يوه بديهي نظام کښې د !!{formula}s په توګه تمثيل کړئ۔
\item په خوندي، «متناهي‌پالو» لارو ثابته کړئ چې دا صوري استنتاجي
  نظامونه سازګار دي۔
\end{enumerate}

د هېلبرټ کار د لومړۍ موخې د پوره کولو پر لور ډېر وړاندې لاړ۔ په 1899
کښې ئې دا کار په خپل نامتو کتاب \emph{د هندسې بنسټونه} کښې د هندسې
دپاره کړے و۔ په ورپسې کلونو کښې هغه او يو شمېر زده‌کوونکو او
همکارانو ئې د رياضياتو پر نورو څانګو کار وکړ، څو هغه څه وکړي چې
هېلبرټ د هندسې دپاره کړي وو۔ هېلبرټ پخپله د حساب او تحليل دپاره
بديهي نظامونه ورکړل۔ څېر مېلو د سټ تيوري بديهي کول ورکړل، چې
فرېنکل، سکولېم، فون نويمان او نورو وغځول۔ د 1920مو کلونو تر نيمايي
دوه داسې لارې وې چې د «ټولو» رياضياتو د بديهي کولو نوم ئې غوښت:
د رسل او وايټ‌هېډ \emph{پرېنسيپيا مېتېماتيکا} او هغه څه چې وروسته
د څېر مېلو--فرېنکل سټ تيوري په نوم وپېژندل شول۔

په 1921 کښې هېلبرټ يوه څېړنيزه پروژه پيل کړه څو د دې نظامونو د
سازګارۍ د ثبوت موخه پوره کړي۔ د هغه څو زده‌کوونکو، په ځانګړي ډول
برنايس، اکرمان، او وروسته ګېنڅن، په دې پروژه کښې ورسره مرسته وکړه۔
د دې موخې بنسټيز نظر دا و چې په رياضياتو کښې د ناسازګارۍ د
!!{derivation} د امکان پوښتنه د نښو د ممکنه لړيو په باب يو
ترکيبي مسئله وګرځوي؛ يعنې د هغو جملو ممکنه لړۍ چې د حساب، تحليل يا
سټ تيورۍ د يوه بديهي نظام له بديهي اصولو څخه د، مثلاً،
$!A \land \lnot !A$ د سم !!{derivation} معيار پوره کوي۔ د نښو د
داسې لړۍ د ناممکن والي ثبوت، څرنګه چې پخپله رياضياتي ثبوت دے، په دې
بديهي نظامونو کښې صوري کېدے شي۔ په بله وينا، داسې يوه جمله~$\OCon$
به وي چې وايي، مثلاً، حساب سازګار دے۔ سربېره پر دې، دا جمله بايد په
اړوندو نظامونو کښې د اثبات وړ وي، په ځانګړي ډول که ثبوت ئې يوازې ډېرې
محدودې، «متناهي‌پالې» لارې غواړي۔

دويمه موخه، چې جوړ شوي بديهي نظامونه به هره رياضياتي پوښتنه حل کړي،
په دوو لارو کره کېدے شي۔ په يوه لار داسې بيانېدے شي: د رياضياتو د
بديهي نظام د ژبې د هرې جملې~$!A$ دپاره يا~$!A$ يا~$\lnot !A$ له
بديهي اصولو څخه د اثبات وړ وي۔ که دا رښتيا وے، داسې جمله به نۀ وه چې
د بديهي اصولو پر بنسټ نه ثابته او نه ردېدے شي، او داسې پوښتنه به نۀ
وه چې بديهي اصول ئې حل نۀ کړي۔ دې خاصيت لرونکی بديهي نظام
\emph{بشپړ} بلل کېږي۔ البته د هرې ټاکلې جملې دپاره لا هم دا معلومول
ستونزمن کار کېدے شي چې له دوو بديلونو څخه کوم يو سم دے۔ خو په اصل
کښې بايد د معلومولو يوه لار وي۔ په حقيقت کښې د هېلبرټ تر څېړنې لاندې
د بديهي اصولو او !!{derivation} نظامونو دپاره بشپړتيا دا لازموله چې
داسې يوه لار شته، که څه هم هېلبرټ په دې نۀ پوهېده۔ د پوښتنې دويم
تعبير دا لا پياوړې غوښتنه ده: داسې ميخانيکي، محاسبوي لار موجوده وي
چې د يوې ورکړل شوې جملې~$!A$ دپاره معلومه کړي له بديهي اصولو څخه
استنتاجېدونکې ده که نۀ۔

په 1931 کښې ګوډل د «ناتکميلۍ دوه قضيې» ثابتې کړې، چې وئې ښودل دا
پروګرام بريالی کېدے نۀ شي۔ د رياضياتو دپاره هېڅ داسې بديهي نظام نشته
چې بشپړ وي؛ په ځانګړي ډول هغه جمله چې د بديهي اصولو سازګاري بيانوي
نه ثابتېدے او نه ردېدے شي۔

دې د هېلبرټ اصلي پروګرام ته وژونکی ګوزار ورکړ۔ خو لکه په رياضياتو
کښې چې ډېر ځله کېږي، د څېړنې په زړه پورې نوې لارې ئې هم پرانيستې۔
که د ټولو رياضياتو يو واحد، هر اړخيز صوري نظام نشته، نو معنا لري چې
لا محدود نظامونه جوړ او وڅېړل شي چې څه پکې ثابتېدے شي۔ همدارنګه معنا
لري چې د دې نظامونو د سازګارۍ د ثابتولو دپاره لږ محدودې ثبوتي لارې
جوړې، او د هغو د سازګارۍ د ثبوت د سختوالي د اندازه کولو لارې وموندل
شي۔ څرنګه چې ګوډل وښودل (نږدې) هر صوري نظام داسې پوښتنې لري چې نۀ
ئې شي حل کولے، نو معنا لري چې د هغو «په زړه پورې» پوښتنو لټون وشي
چې يو ورکړل شوے صوري نظام ئې نۀ شي حل کولے، او معلومه شي چې د حل
دپاره ئې صوري نظام بايد څومره پياوړے وي۔ منطقپوهان تر نن ورځې دا
پوښتنې د رياضياتو په يوه نوې څانګه، د ثبوتونو په تيورۍ، کښې څېړي۔

\end{document}
